% !TeX program = xelatex
% 支撑材料：AI 工具使用详情（依据《全国大学生数学建模竞赛人工智能工具使用规定（2026年试行）》第四条）
\documentclass[12pt,a4paper]{ctexart}
\usepackage[a4paper,top=2.6cm,bottom=2.6cm,left=2.6cm,right=2.6cm]{geometry}
\usepackage{booktabs,array,enumitem,hyperref}
\usepackage{listings}
\ctexset{
  section = {format = \heiti\zihao{4}, name = {,、}, number = \chinese{section}, aftername = {}},
  subsection = {format = \heiti\zihao{-4}, aftername = {}}
}
\linespread{1.25}
\setlength{\parindent}{2em}
\lstset{basicstyle=\ttfamily\footnotesize, breaklines=true, frame=single,
        extendedchars=false, columns=flexible, showstringspaces=false}

\begin{document}

\begin{center}
  {\heiti\zihao{3} AI 工具使用详情}\\[4pt]
  {\zihao{-4} （2026 年高教社杯全国大学生数学建模竞赛\quad C 题\quad 支撑材料）}
\end{center}

\section*{一、所用 AI 工具的名称、版本或型号}

\begin{table}[h]
  \centering
  \small
  \begin{tabular}{p{3.6cm}p{4.2cm}p{7.6cm}}
    \toprule[1.5pt]
    类别 & 名称与版本 & 用途概要 \\
    \midrule[1pt]
    AI 智能体平台 & DeepSeek Harness 桌面端（DSH Desktop） & 提供本次建模的对话式编码智能体环境 \\
    对话模型 & deepseek-v4-pro 与 deepseek-flash（会话中发生一次模型切换） & 思路讨论、代码生成、文档撰写 \\
    建模插件 & UltraMath 数学建模多智能体预设（含 33 篇模型方法库） & 读取题目、方法选型索引、阶段门禁检查 \\
    编程环境 & Python 3.12.7、numpy 1.26.4、pandas 2.2.2、scipy 1.13.1（HiGHS）、matplotlib 3.9、openpyxl 3.1 & 线性规划求解、数据处理与绘图 \\
    排版工具 & XeLaTeX（TeX Live 2026） & 论文与本文档的排版编译 \\
    \bottomrule[1.5pt]
  \end{tabular}
\end{table}

\section*{二、具体使用目的和环节}

AI 工具仅用于\textbf{辅助性}工作，模型建立的核心判断、结果取舍与结论均由参赛队员作出。具体环节如下：

\begin{enumerate}[label=(\arabic*),leftmargin=2.2em,itemsep=2pt]
  \item \textbf{题目与数据理解}：辅助阅读题目与附件结构，编写数据探索脚本（统计附件1--4 的规模、量纲、取值范围、月度与日内规律），并由队员核对统计口径。
  \item \textbf{建模框架与算法选型讨论}：就"确定性线性规划、动态规划、两阶段随机规划、滚动时域优化、鲁棒优化"等候选方法作适用性比较，形成"以线性规划为统一内核"的技术路线；该路线由队员确认后实施。
  \item \textbf{程序编写与调试}：生成线性规划建模与求解代码（变量与约束装配、HiGHS 调用）、预测器实现（滑动平均、指数加权平均）、滚动调整与费用结算代码、结果文件按附件5 模板写出，以及独立验算代码。
  \item \textbf{图表生成}：编写绘图脚本，从求解结果文件直接生成论文全部插图。
  \item \textbf{论文撰写与排版}：辅助组织章节结构、撰写文字初稿、排版 LaTeX 源码与编译排错；语言由队员逐段审阅修改。
  \item \textbf{独立验算}：生成"母线侧变量布局的独立线性规划"与"储电量网格动态规划"两套交叉验算代码。
\end{enumerate}

\section*{三、主要提示方式与使用过程说明（典型交互示例）}

使用方式为"分阶段任务提示 $+$ 结果核验 $+$ 纠偏迭代"，即每次给出明确的目标、输入数据范围与验收标准，并规定"不得使用未来数据""结果需给出可复核的统计量"等约束。以下为四个典型交互示例。

\textbf{示例 1（数据口径确认）}

\noindent\fbox{\parbox{0.97\linewidth}{\small
\textbf{提问}：附件2 与附件4 的时间列是 \texttt{00:10,\,00:20,\dots,0:00+1}，而附件5 结果模板的时段写作 \texttt{0:10-0:20,\dots}。请判断二者如何对应，并说明依据；不要改动任何数据文件。\\
\textbf{回答要点}：指出附件时标为"区间结束时刻"（\texttt{00:10} 表示 $[0{:}00,0{:}10]$，\texttt{0:00+1} 表示 $[23{:}50,24{:}00]$），模板按时段起点命名，二者相差一个时段的命名偏移；建议全文统一取前者，并给出"144 个标签恰好覆盖 $[0,24]$ 小时"的校验。\\
\textbf{队员处理}：采纳该口径，并要求在代码中与论文口径声明中同时写清，避免图表与数据错位。
}}

\textbf{示例 2（模型选择与机理推导）}

\noindent\fbox{\parbox{0.97\linewidth}{\small
\textbf{提问}：问题2 中缺电按 5 倍电价紧急购电，计划应留多少安全裕度？请给出理论依据，并与数值标定相互验证。\\
\textbf{回答要点}：将其归结为报童问题，缺货成本 $c_u=5p_t-p_t=4p_t$，多备成本 $c_o=p_t(1-\eta_c\eta_d)=0.19p_t$（储能有余量）或 $p_t$（无余量），故最优覆盖分位点为 $c_u/(c_u+c_o)\in[0.80,0.96]$；并提示"$\eta$、$p_t$ 同比例缩放时分位点不变、绝对代价随时段变化"，即储备应优先分配高价时段。\\
\textbf{队员处理}：核验推导后，用 1 月数据网格搜索标定得到 $(\alpha,\beta)=(0.15,0)$，与理论区间一致；将该推导写入论文 \S4.4。
}}

\textbf{示例 3（结论反直觉时的排查）}

\noindent\fbox{\parbox{0.97\linewidth}{\small
\textbf{提问}：问题3 出现"增加预报更新反而使总费用上升"的结果，请判断是数据特征还是程序缺陷，并给出可验证的排查步骤。\\
\textbf{回答要点}：先列三条可检验判据——(i) 分解费用为计划费、违约费、紧急费，比较各策略；(ii) 计算"同一时间窗内"不同发布时刻的预报精度，检验后续预报是否真的更准；(iii) 统计相邻发布时刻的预报修订幅度（同一时钟小时对齐）。据此定位到两点：违约费增量大于紧急费下降量，且 19:00--24:00 的四个发布时刻预报完全相同（夜间光伏为零，无新信息）。\\
\textbf{队员处理}：按该排查路径复算，确认结论成立；同时要求把"预报对齐口径"写入假设并在检验章节给出敏感性分析。
}}

\textbf{示例 4（结果复核）}

\noindent\fbox{\parbox{0.97\linewidth}{\small
\textbf{提问}：请对问题1 的最优值给出与主程序相互独立的第二种算法，并要求说明二者偏差的合理范围。\\
\textbf{回答要点}：给出两条独立路线——把变量改为"母线侧充电量/放电量"的等价线性规划（变量个数与约束矩阵均不同），以及对储电量按 20 kWh 网格离散的动态规划；并说明动态规划是原问题的受限离散解，其最优值应\textbf{不低于}线性规划最优值，偏差只来自离散化。\\
\textbf{队员处理}：运行两套独立代码，得到 $35101.6$ 元与 $35198.4$ 元（间隙 $0.276\%$），符合预期后写入论文 \S7.1。
}}

\section*{四、对 AI 输出的采纳、人工修改和核验情况}

\textbf{（1）采纳与修改。} 采纳的内容包括：候选方法清单与选型比较、程序的框架与大部分实现、图表绘制脚本、文字初稿与 LaTeX 排版源码。人工修改包括：确定最终技术路线（统一线性规划内核）；确定并统一全部口径假设（时间区间口径、效率与功率口径、终端电量口径、附件3 预报对齐口径、违约费基准口径、问题4 结算价口径）；改写摘要与各章关键论述；调整图表的内容与呈现方式。

\textbf{（2）核验措施。} 参赛队对 AI 参与完成的内容逐项核验，主要措施有四条：

\begin{enumerate}[label=(\arabic*),leftmargin=2.2em,itemsep=2pt]
  \item \textbf{推导复核}：论文中全部公式与命题（购电量分解式、充放电互斥性、阈值策略、报童分位点、对称罚金等价式、中位数准则、更新价值分解、价格相关性下分位点上移）均要求给出推导或证明过程，由队员逐条复核其前提与量纲。
  \item \textbf{独立验算}：另写一套不共用求解代码的程序，用\textbf{不同的算法}（母线侧布局 LP、储电量网格 DP）与\textbf{独立结算公式}复核核心结果；共 40 项检查，全部通过。
  \item \textbf{结果文件核对}：直接读取 5 个结果文件，检验 144 个时段值的完整性、"全天合计 $=$ 行和"、逐日 $0:00$ 与 $24:00$ 储电量、"充放电量 $-$ 储电量"自洽性以及储电量与功率的界约束。
  \item \textbf{数字溯源}：论文中的每一个数值均标注来源（结果文件或统计脚本输出），并由队员逐项比对，图表全部由脚本直接生成、未作人工修饰。
\end{enumerate}

\textbf{（3）AI 输出的纠错记录。} 核验过程中发现并纠正了 AI 生成内容中的三类错误，均经人工核验后修正：

\begin{enumerate}[label=(\roman*),leftmargin=2.4em,itemsep=2pt]
  \item \textbf{预报对齐口径错误}：最初按"当日第 $k$ 小时"理解附件3 的"预报 $k$ 小时"，导致 6:00、12:00、18:00 发布的预报与时段错位 $6$--$18$ 小时，问题3 的结论出现反转；经"同一时间窗精度比较"检验后改为"发布时刻之后第 $k$ 小时"，结论随之修正。
  \item \textbf{结果费用列索引错误}：问题4 结果文件的"全天购电费"列一度按日期行号错位取用电价，使 12 月等日期数值偏低约 $45\%$；由独立复算发现并修正。
  \item \textbf{数据读取错误}：附件4 读取函数在默认参数下返回了字典而非数据表，导致问题4 首次运行失败；修正后重跑。
\end{enumerate}

上述记录表明：AI 工具在本项目中承担的是"助手"角色，其输出均经过参赛队员的复核、复算与纠错，模型的核心建模与分析由参赛队主导完成。

\end{document}
